% Part: normal-modal-logic
% Chapter: completeness
% Section: completeness-K

\documentclass[../../../include/open-logic-section]{subfiles}

\begin{document}

\olfileid{nml}{com}{cmk}

\olsection{\Log{K}కు నిర్ణాయకత్వం, సంపూర్ణత}

ఇప్పుడు కానానికల్ నమూనాతో సంపూర్ణతను స్థాపించడానికి
సిద్ధంగా ఉన్నాం. $\Sigma$కు చెందిన కానానికల్ నమూనాలో
సత్యమయ్యే \tetoken{సూత్రాలు}{!!{formula}s},
$\Sigma$ నుంచి \tetoken{వ్యుత్పాద్యమైనవే}{!!{derivable}}
అనే వాస్తవం వల్ల సంపూర్ణత అనుసరిస్తుంది. ఈ ధర్మమున్న
నమూనాలు~$\Sigma$ను \emph{నిర్ణయిస్తాయి} అంటాం.

\begin{defn}
  ప్రతి \tetoken{సూత్రం}{!!{formula}s}~$!A$కు
  $\mSat{M}{!A}$, $\Sigma \Proves !A$ రెండూ
  సమానార్థకమైనప్పుడే నమూనా $\mModel{M}$ నార్మల్
  మోడల్ తర్కం~$\Sigma$ను \emph{నిర్ణయిస్తుంది}.
\end{defn}

\begin{thm}[నిర్ణాయకత్వం]\ollabel{thm:determination}
   $\mSat{M^\Sigma}{!A}$ అయితే, అప్పుడు మాత్రమే
   $\Sigma \Proves !A$.
\end{thm}

\begin{proof}
  $\mSat{M^\Sigma}{!A}$ అయితే, ప్రతి సంపూర్ణ
  $\Sigma$-అవిరుద్ధ $\Delta$కు
  $\mSat{M^\Sigma}{!A}[\Delta]$. కనుక సత్య ఉపసిద్ధాంతం వల్ల
  ప్రతి సంపూర్ణ $\Sigma$-అవిరుద్ధ $\Delta$లో $!A \in
  \Delta$. అందువల్ల \olref[lin]{cor:provability-characterization}
  ప్రకారం ($\Gamma = \emptyset$ తీసుకుంటే)
  $\Sigma \Proves !A$.

  విపరీతంగా, $\Sigma \Proves !A$ అయితే
  \olref[ccs]{prop:ccs-properties}\olref[ccs]{prop:ccs-closed}
  ప్రకారం ప్రతి సంపూర్ణ $\Sigma$-అవిరుద్ధ $\Delta$లో
  $!A$ ఉంటుంది. కనుక సత్య ఉపసిద్ధాంతం వల్ల ప్రతి~$\Delta \in
  W^\Sigma$కు $\mSat{M^\Sigma}{!A}[\Delta]$; అంటే,
  $\mSat{M^\Sigma}{!A}$.
\end{proof}

\Log{K}కు చెందిన కానానికల్ నమూనా~\Log{K}ను నిర్ణయిస్తుంది
కనుక, \Log{K} సంపూర్ణత వెంటనే అనుబంధ ఫలితంగా వస్తుంది:

\begin{cor}\ollabel{cor:Kcomplete}
  ప్రాథమిక మోడల్ తర్కం \Log{K} అన్ని నమూనాల వర్గానికి
  సంబంధించి సంపూర్ణం; అంటే, $\Entails !A$ అయితే
  $\Log{K} \Proves !A$.
\end{cor}

\begin{proof}
  విపర్యయంగా, $\Log{K} \Proves/ !A$ అయితే నిర్ణాయకత్వం
  వల్ల $\mSat/{M^{\Log{K}}}{!A}$; కనుక $!A$ సర్వత్రా
  చెల్లుబాటు కాదు.
\end{proof}

ఒక నమూనా వర్గానికి సంబంధించి వ్యవస్థ $\Sigma$ సంపూర్ణతను
సాధారణంగా నిరూపించాలంటే---ఉదాహరణకు, స్వావర్తన, సౌష్ఠవ,
సంక్రమణీయ నమూనాల వర్గానికి సంబంధించి \Log{KTB4}
సంపూర్ణతను---నిర్ణాయకత్వం మాత్రమే సరిపోదు. వ్యవస్థ
$\Sigma$కు చెందిన కానానికల్ నమూనా ఆ వర్గంలో ఉందని కూడా
చూపాలి. ఇది కానానికల్ నమూనా నిర్మాణం నుంచే స్పష్టంగా
అనుసరించదు---ఎప్పుడూ నిజం కూడా కాదు!

\end{document}
